\chapter{Input and Simulation Setup}
\label{app:input}

This appendix illustrates how the elements defined in Chapter~\ref{chapter:proposed_model} can be represented when preparing a LODUS simulation. A simulation setup comprises: 
\begin{itemize} 
    \item \textit{i)} environment data describing regions, points of interest, their attributes, and dependency requirements (Section~\ref{app:input_environment}); 
    \item \textit{ii)} population data describing the initial blobs and their sampled and traceable characteristics (Section~\ref{app:input_population}); 
    \item \textit{iii)} routine data describing local and global scheduled actions (Section~\ref{app:input_routine}); 
    \item \textit{iv)} an experiment configuration defining the simulation horizon, selected plugins, and their parameters (Section~\ref{app:experiment_configuration}); and 
    \item \textit{v)} the simulation procedure that combines these elements and advances the simulation over time (Section~\ref{app:simulator}). 
\end{itemize} 

The examples are deliberately small and illustrate how the abstractions of the proposed model can be instantiated for a simulation. The specific data sources and attributes depend on the application domain, while the representation of regions, points of interest, blobs, templates, routines, and actions follows the definitions introduced in Chapter~\ref{chapter:proposed_model}.


\section{Environment Input Files}
\label{app:input_environment}

The environment input represents the regions $R$ and points of interest $P$ defined in Section~\ref{sec:environment}. A region is identified by a name, a reference coordinate, and a collection of points of interest. Each point of interest has an identifier, a semantic type corresponding to $\tau_i$, a coordinate, and an optional collection of attributes corresponding to $\bm{\alpha}_i$.

Listing~\ref{lst:environment_description} illustrates an environment containing a home, a school, a pharmacy, and infrastructure points of interest. The attributes associated with each point of interest represent properties that may be relevant to the corresponding experiment. For example, a capacity may be used in a service scenario, while a water-level threshold may be used in a flood scenario. Their interpretation depends on the behavior defined for the experiment.

\begin{lstlisting}[language={Python},
caption={Example of an environment input file with typed POIs and attributes.},
captionpos=b, label={lst:environment_description}, basicstyle=\scriptsize]
{
  "regions": [
    {
      "name": "ExampleRegion",
      "lng_lat": [-51.19757, -30.15647],
      "points_of_interest": [
        {
          "poi_type": "home",
          "unique_name": "home_0",
          "lng_lat": [-51.19805, -30.15633],
          "attributes": {"external_id": "H-001"}
        },
        {
          "poi_type": "school",
          "unique_name": "school_0",
          "lng_lat": [-51.19750, -30.15621],
          "attributes": {"capacity": 120}
        },
        {
          "poi_type": "pharmacy",
          "unique_name": "pharmacy_0",
          "lng_lat": [-51.19820, -30.15655],
          "attributes": {"capacity": 40, "water_level": 3.5}
        },
        {
          "poi_type": "power_source",
          "unique_name": "power_source_0",
          "lng_lat": [-51.19910, -30.15590]
        },
        {
          "poi_type": "water_source",
          "unique_name": "water_source_0",
          "lng_lat": [-51.19940, -30.15560]
        },
        {
          "poi_type": "water_source",
          "unique_name": "water_source_1",
          "lng_lat": [-51.19680, -30.15570]
        }
      ]
    }
  ]
}
\end{lstlisting}

The direct disabling conditions $\bm{C}_i(t)$ represent simulation state and therefore do not need to be specified as permanent availability values in the environment description. Such conditions may be introduced or removed during the simulation, for example when a point of interest closes or becomes exposed to flooding.

Dependency requirements $D_i$ may also be supplied as part of the environment configuration. As defined in Section~\ref{sec:environment}, each requirement is represented by a pair $(k_{ir},\bm{G}_{ir})$, where $\bm{G}_{ir}$ contains a set of prerequisite points of interest and $k_{ir}$ specifies the minimum number that must be available.

For example, the pharmacy introduced in Listing~\ref{lst:environment_description} may require an available power source and at least one of two available water sources. These requirements correspond to

\[ D_p = \left\{ (1,\{P_e\}), (1,\{P_{w_1},P_{w_2}\}) \right\}. \] 

The first requirement represents a mandatory dependency on $P_e$, while the second allows either $P_{w_1}$ or $P_{w_2}$ to satisfy the water dependency. All requirements contained in $D_p$ must be satisfied for the pharmacy to be available. 

Changes in availability do not directly modify the population occupying a point of interest. Any cancellation, waiting, movement, or redirection of that population is expressed through the actions and domain-specific behavior defined for the simulation.

\begin{lstlisting}[language={Python},
caption={Example of all-of and minimum-of POI dependency rules.},
captionpos=b, label={lst:dependency_description}, basicstyle=\scriptsize]
{
  "dependency_rules": {
    "pharmacy_infrastructure": {
      "all_of": [
        "ExampleRegion//power_source_0"
      ],
      "min_of": [
        {
          "minimum_enabled": 1,
          "nodes": [
            "ExampleRegion//water_source_0",
            "ExampleRegion//water_source_1"
          ]
        }
      ]
    }
  },
  "node_dependencies": {
    "ExampleRegion//pharmacy_0": ["pharmacy_infrastructure"]
  }
}
\end{lstlisting}

The simulator recomputes effective availability after a direct state change
and propagates dependency effects until the states stabilize. Unavailability
does not itself move, delete, or queue a blob. The handler consuming an action
must explicitly decide whether to suppress, retain, redirect, or otherwise
transform the affected demand.


\section{Population Input Files}
\label{app:input_population}

The population input describes the initial blobs $B=(TP,\bm{Tc},\bm{Sc})$ defined in Section~\ref{sec:population}. Initial values may be specified for traceable characteristics that are shared by all individuals represented by a blob. A blob may then provide values that differ from these defaults when necessary. 

The population input also defines the categories used by sampled characteristics. Within each blob, the category counts associated with each sampled characteristic sum to $TP$. As discussed in Section~\ref{sec:population}, these values describe marginal distributions and do not establish joint relationships between different sampled characteristics. 

Additional traceable characteristics required by a particular experiment may be introduced through the model extensions $E_{\Pi}$ of the selected plugins. Once introduced, these characteristics follow the same representation and participate in blob selection, splitting, and merging. 

The initial population associates points of interest with their respective collections $\bm{B}_{P_i}$. Points of interest with no initial population may be omitted. In Listing~\ref{lst:population_description}, blob $B_i$ contains 100 unvaccinated individuals, while blob $B_j$ contains 20 vaccinated children. For each blob, the age and occupation distributions independently sum to its total population.

\begin{lstlisting}[language={Python},
caption={Example of a population input file.},
captionpos=b, label={lst:population_description}, basicstyle=\scriptsize]
{
  "default_traceable_characteristics": {
    "vaccination_status": "unvaccinated"
  },
  "sampled_characteristics_categories": {
    "age": ["children", "youngs", "adults", "elders"],
    "occupation": ["student", "worker", "other"]
  },
  "initial_population": {
    "ExampleRegion//home_0": [
      {
        "total_population": 100,
        "traceable_characteristics": {},
        "sampled_characteristics": {
          "age": {
            "children": 20,
            "youngs": 10,
            "adults": 50,
            "elders": 20
          },
          "occupation": {
            "student": 30,
            "worker": 50,
            "other": 20
          }
        }
      },
      {
        "total_population": 20,
        "traceable_characteristics": {
          "vaccination_status": "vaccinated"
        },
        "sampled_characteristics": {
          "age": {
            "children": 20,
            "youngs": 0,
            "adults": 0,
            "elders": 0
          },
          "occupation": {
            "student": 20,
            "worker": 0,
            "other": 0
          }
        }
      }
    ]
  }
}
\end{lstlisting}

All traceable characteristics considered by the simulation participate in blob compatibility. Thus, blobs at the same point of interest can merge only when their traceable-characteristic values agree. When compatible blobs are merged, the corresponding sampled-category counts are added.


\section{Routine Input Files}
\label{app:input_routine}

Routine inputs represent the scheduled actions $(t,A)$ that compose the local routines $\bm{Ro}_{P_i}$ and the global routine $\bm{GRo}$ described in Section~\ref{sec:routines}. Each scheduled entry specifies a cycle step $t$ and an action 

\[ A_i=(\kappa_i,T_i,\bm{V}_i), \] 

where $\kappa_i$ identifies the action type, $T_i$ specifies the eligible population, and $\bm{V}_i$ contains any additional parameters required by the action. 

In the input representation, the population-template component describes the traceable- and sampled-characteristic predicates associated with $T_i$, while the action values describe the parameters represented by $\bm{V}_i$. The requested quantity corresponds to $T_Q$ in the population-template definition. 

Listing~\ref{lst:routine_description} illustrates scheduled actions in local and global routines. A local action is associated with a particular point of interest, whereas a global action is scheduled independently of a single location. When multiple actions are scheduled for the same cycle step, their ordering is preserved because each action observes the simulation state resulting from the actions applied before it.

\begin{lstlisting}[language={Python},
caption={Example of ordered local and global routine inputs.},
captionpos=b, label={lst:routine_description}, basicstyle=\scriptsize]
{
  "global_routine": [
    {
      "cycle_step": [10],
      "execution_scope": "once",
      "action": {
        "name": "return_pharmacy_visitors",
        "type": "example_action",
        "population_template": {
          "traceable_characteristics": {
            "vaccination_status": "unvaccinated"
          },
          "sampled_characteristics": {}
        },
        "values": {
          "region": "ExampleRegion",
          "node": "pharmacy_0",
          "quantity": 20
        }
      }
    }
  ],
  "routines": {
    "ExampleRegion//home_0": [
      {
        "cycle_step": 8,
        "action": {
          "name": "send_students_to_pharmacy",
          "type": "move_population",
          "population_template": {
            "traceable_characteristics": {
              "vaccination_status": "unvaccinated"
            },
            "sampled_characteristics": {
              "age": ["children"],
              "occupation": ["student"]
            }
          },
          "values": {
            "origin_region": "ExampleRegion",
            "origin_node": "home_0",
            "destination_region": "ExampleRegion",
            "destination_node": "pharmacy_0",
            "quantity": 20
          }
        }
      }
    ]
  }
}
\end{lstlisting}


\section{Experiment Configuration}
\label{app:experiment_configuration}

The experiment configuration combines the structural inputs and the domain-specific extensions required by a simulation. It identifies the environment, population, and routine inputs, defines the temporal parameters of the experiment, and specifies the plugins and their corresponding parameters. 

The temporal configuration includes the number of cycles and their length $d$. Consequently, a simulation containing $c$ cycles contains $cd$ global simulation steps. At global step $s$, the corresponding cycle step is obtained from 

\[ t=s\bmod d. \] 

For each selected plugin $\Pi$, the experiment configuration may also provide the parameters represented by $\theta_{\Pi}$. These values determine the particular configuration used in the experiment without changing the core environment, population, routine, or action definitions. 

Listing~\ref{lst:experiment_configuration} illustrates how the input files, temporal parameters, and plugin-specific parameters can be combined in one experiment description.

\begin{lstlisting}[language={Python},
caption={Example of an experiment configuration file.},
captionpos=b, label={lst:experiment_configuration}, basicstyle=\scriptsize]
{
  "envgraph_inputs_files": {
    "environment_file": "Example-Environment.json",
    "population_file": "Example-Population.json",
    "routine_file": "Example-Routine.json"
  },
  "simulation_parameters": {
    "total_cycles": 3,
    "cycle_length": 24,
    "random_seed": 0
  },
  "node_dependency_data_plugin": {
    "dependency_files": ["Example-Dependencies.json"]
  },
  "example_plugin": {
    "destination_node": "home_0",
    "default_quantity": 20
  }
}
\end{lstlisting}


\section{Plugin Contributions}
\label{app:plugin}

Plugins provide the domain-specific extensions described in Section~\ref{sec:extension_mechanism}. Their contributions are summarized by 

\[ \Pi = \left( \theta_{\Pi}, E_{\Pi}, R_{\Pi}, F_{\Pi}, H_{\Pi} \right), \] 

where $\theta_{\Pi}$ contains experiment-specific parameters, $E_{\Pi}$ contains extensions to the model representation, $R_{\Pi}$ contains routines or scheduled actions, $F_{\Pi}$ contains additional action types and their semantics, and $H_{\Pi}$ contains other simulation-step updates, observations, data access, or output recording associated with the plugin. 

A plugin does not need to contribute all five components. For example, a plugin used to represent an environmental disruption may provide parameters through $\theta_{\Pi}$ and update point-of-interest availability through $H_{\Pi}$ without introducing a new population characteristic. Conversely, an epidemic plugin may introduce an infection-state characteristic through $E_{\Pi}$ and additional actions through $F_{\Pi}$. 

The experiment configuration determines which plugins are included and the values of their parameters. Once included, plugin-defined actions retain the same representation $A_i=(\kappa_i,T_i,\bm{V}_i)$ used by the core actions. Consequently, domain-specific behavior extends the model without introducing a separate representation for population selection or state changes. 

This representation also provides a common way of describing the plugins used in the experimental scenarios. Only the components relevant to a particular plugin need to be reported.

\begin{comment}
    

\begin{lstlisting}[language={Python},
caption={Example of a configured composite action plugin.},
captionpos=b, label={lst:plugin_script}, basicstyle=\scriptsize]
from core.plugin import ActionPlugin
from core.routine import Action


class ExamplePlugin(ActionPlugin):

    def load_plugin(self, simulation):
        self.simulation = simulation
        self.graph = simulation.env_graph
        self.config = simulation.experiment_config.get(
            "example_plugin", {}
        )

        # Delta Sigma: extend every blob with a default value.
        
        self.graph.add_blobs_traceable_property(
            "example_state", "idle"
        )

        # F: register a composite action handler.
        simulation.add_action_type_to_function(
            "example_action", self.consume_example_action, False
        )

    def consume_example_action(
        self, pop_template, values, cycle_step, simulation_step
    ):
        destination = self.config.get(
            "destination_node", "home_0"
        )
        quantity = values.get(
            "quantity", self.config.get("default_quantity", 20)
        )

        move_values = {
            "origin_region": values["region"],
            "origin_node": values["node"],
            "destination_region": values["region"],
            "destination_node": destination,
            "quantity": quantity
        }
        return [Action("move_population", pop_template, move_values)]

    def update_time_step(self, cycle_step, simulation_step):
        pass

    def unload_plugin(self):
        pass
\end{lstlisting}
\end{comment}

\section{Simulation Procedure}
\label{app:simulator}

A simulation combines the configured environment, initial population, routines, and selected plugins and advances them over a sequence of discrete simulation steps. 

First, the regions and points of interest are created from the environment inputs and their initial populations are represented by the corresponding blobs. The local routines $\bm{Ro}_{P_i}$ and global routine $\bm{GRo}$ are then associated with the simulation, together with the extensions and parameters contributed by the selected plugins. 

For each global simulation step $s$, the corresponding cycle step is 

\[ t=s\bmod d. \] 

The core action-processing procedure then follows the sequence defined in Section~\ref{sec:routines}: 

\begin{enumerate} 
\item collect the actions scheduled for cycle step $t$ by $\bm{GRo}$ and the local routines $\bm{Ro}_{P_i}$; 
\item apply these actions sequentially in their defined order; and 
\item merge compatible blobs located at the same point of interest after the scheduled actions have been applied. 
\end{enumerate} 

Because actions are applied sequentially, each action operates on the state produced by the actions preceding it at the same cycle step. Population templates are therefore evaluated against the current simulation state when their corresponding actions are applied. 

Plugins may additionally contribute simulation-step updates, observations, or other domain-specific behavior through $H_{\Pi}$. These contributions augment the procedure above without changing the common representations of points of interest, blobs, templates, routines, and actions. 

The process is repeated until the configured number of simulation steps has been completed.

\begin{comment}
    

\begin{lstlisting}[language={Python},
caption={Example of a simulator script using the revised lifecycle.},
captionpos=b, label={lst:simulator}, basicstyle=\scriptsize]
from plugins.data.node_dependency_data_plugin import (
    NodeDependencyDataPlugin
)
from plugins.time_actions.move_population_plugin import (
    MovePopulationPlugin
)
from util.data_parse import generate_lodus_simulation

from example_plugin import ExamplePlugin


# Parse experiments/ExampleExperiment.json and its input files.
simulation = generate_lodus_simulation("ExampleExperiment")

# Load the extension set required by the configured actions.
simulation.load_plugin(NodeDependencyDataPlugin())
simulation.load_plugin(MovePopulationPlugin())
simulation.load_plugin(ExamplePlugin())

# Advance all cycles. The simulation object manages t and s.
number_of_steps = (
    simulation.time_status.total_cycles
    * simulation.time_status.cycle_length
)
for _ in range(number_of_steps):
    simulation.update_time_step()
\end{lstlisting}

\end{comment}